\section{Related Works}

Helped by their accessibility, there exist multiple prior works on ``legal NLI'' for case and statute laws.
One of the subtasks in COLIEE-2020 shared task \cite{rabelo_coliee_nodate} was, given a court decision Q and relevant cases, to extract relevant paragraphs from the cases and to classify whether those paragraphs entail ``Q'' or ``not Q''.
\citet{holzenberger_dataset_2020} introduced a dataset for predicting an entailment relationship between a statement and a statute excerpt.
While they are both ``legal'' and ``NLI'', statutes and contracts exhibit different characteristics including the fact that statutes/cases tend to be written in consistent vocabulary and styles.
Moreover, there only exists a single right answer for a hypothesis in case/statute law NLI, whereas a hypothesis can be entailed by or contradicting to each contract in our task; i.e., hypotheses and documents have one-to-one relationships in case/statute law NLI, but they have many-to-many relationships in our task.

As discussed in \cref{sec:introduction}, our task has practical and scientific significance compared to information extraction for contracts \cite{leivaditi_benchmark_2020,hendrycks_cuad_2021}.
We showed in our experiments that the NLI part of our task is much more challenging than the evidence identification task.
% Unlike \cite{hendrycks_cuad_2021}, which were annotated by multiple professionals, our dataset might not be as faithful to an existing law as theirs due to less involvement of law professionals.
% We argue that annotation based on guideline .. and our dataset should therefore be good enough as a proxy for studying legal data.
Furthermore, we gave observations to linguistic characteristics of our dataset that are lacking in these prior works.

\citet{lippi_claudette_2019} presented a dataset where certain types of contract clauses are identified and annotated with ``clearly fair'', ``potentially unfair'' or ``clearly unfair''.
While the format of the task input and output is quite similar, our task requires reasoning over a much diverse set of hypotheses than just fair or unfair.
Similarly, fact extraction and claim verification tasks \cite{thorne_fever_2018,jiang_hover_2020}, where the task is to extract facts from Wikipedia articles and to classify whether the claim is entailed by the facts, have similar input and output formats.
Such claims and our hypotheses are quite different in nature and working on contracts poses unique challenges as discussed in \cref{sec:discussion-observations}.
